% ============================================================================
%  Preamble: packages, theorem environments, and notation macros.
%  Shared by main.tex (% ============================================================================
%  Preamble: packages, theorem environments, and notation macros.
%  Shared by main.tex (% ============================================================================
%  Preamble: packages, theorem environments, and notation macros.
%  Shared by main.tex (% ============================================================================
%  Preamble: packages, theorem environments, and notation macros.
%  Shared by main.tex (\input{preamble}). Edit document-wide formatting here.
% ============================================================================

% --- Packages ---
\usepackage[utf8]{inputenc}
\usepackage[T1]{fontenc}
\usepackage{lmodern}
\usepackage[margin=1in]{geometry}
\usepackage{amsmath,amssymb,amsthm}
\usepackage{booktabs}
\usepackage{graphicx}
\usepackage{enumitem}
\usepackage{xcolor}
\usepackage[colorlinks=true,linkcolor=blue!70!black,citecolor=green!50!black,urlcolor=blue!70!black]{hyperref}
\usepackage[round]{natbib}
\usepackage{tikz}
\usetikzlibrary{arrows.meta,decorations.pathreplacing,calc,positioning}
\usepackage[most]{tcolorbox}

% --- Theorem environments ---
\newtheorem{theorem}{Theorem}
\newtheorem{proposition}[theorem]{Proposition}
\newtheorem{definition}[theorem]{Definition}
\newtheorem{remark}[theorem]{Remark}
\newtheorem{lemma}[theorem]{Lemma}
\newtheorem{corollary}[theorem]{Corollary}
\newtheorem{conjecture}[theorem]{Conjecture}
\numberwithin{theorem}{section}

% --- Notation ---
\newcommand{\R}{\mathbb{R}}
\newcommand{\M}{\mathrm{M}}
\newcommand{\Mfx}{M(W{,}f)(x)}
\newcommand{\Mfxp}{M(W{,}f)(x')}
\newcommand{\Weff}{W_{\text{eff}}}
\newcommand{\beff}{b_{\text{eff}}}
\newcommand{\KM}{\textsc{KM}}

% Placeholder macro for empirical numbers wired in by wire_paper_results.py.
% Renders the argument verbatim in the PDF until substituted.
\newcommand{\TODO}[1]{\textcolor{red}{\textbf{[TODO:}\ #1\textbf{]}}}

% ===== REVIEW-ONLY WARNING BOXES (remove every \WARNING{...} before submission) =====
% Red-framed [WARNING] box with a detailed explanation, visible in the compiled PDF.
\newcommand{\WARNING}[1]{%
  \par\smallskip\noindent
  \fcolorbox{red}{red!6}{%
    \begin{minipage}{0.93\linewidth}%
      \small\textcolor{red}{\textbf{[WARNING --- review note; remove before submission]}}\par\smallskip
      #1%
    \end{minipage}}\par\smallskip}

% ===== INLINE CLARIFICATIONS (added 2026-09-07) =====
% Blue inline text answering reader questions: definitions of notation and
% terminology given at the point of first use, rather than in a glossary.
%   \clarify{...}     inline blue prose (a sentence or a short paragraph)
%   \clarifyterm{t}{d} bolds the term being defined, then the blue definition
% Both are plain \textcolor wrappers -- they add no vertical space and can be
% used inside theorem statements, captions, and tables.
\definecolor{clarifyblue}{RGB}{0,70,190}
\newcommand{\clarify}[1]{\textcolor{clarifyblue}{#1}}
\newcommand{\clarifyterm}[2]{\textcolor{clarifyblue}{\emph{#1}: #2}}
% Display-style clarification: an indented blue block for multi-sentence notes.
\newenvironment{clarifybox}
  {\par\smallskip\begingroup\color{clarifyblue}\small\leftskip=1.2em\rightskip=1.2em\noindent\ignorespaces}
  {\par\endgroup\smallskip}

% ===== NETWORK FUNCTION: Psi(W,f) (notation fix, 2026-09-07) =====
% The paper previously used f for BOTH the activation function (as in the
% knowledge matrix M(W,f), where f IS the activation) and the network function
% (as in f(theta,x)). That collision is resolved by writing the network function
% in the quiver-representation notation of Armenta & Jodoin: Psi(W,f), the
% function realised by weight collection W and activation f. From here on
%   f      = the activation function
%   Psi(W,f) = the network function it realises
\newcommand{\Net}{\Psi(W\!,f)}          % the network function itself
\newcommand{\Netx}{\Psi(W\!,f)(x)}      % evaluated at x
\newcommand{\Nety}{\Psi(W\!,f)(y)}
\newcommand{\Netxp}{\Psi(W\!,f)(x')}
\newcommand{\Netarg}[1]{\Psi(W\!,f)(#1)}

% ===== KNOWLEDGE MATRIX indexed by (W,f), not by theta (2026-09-07) =====
% Companion to the \Net macros above. The paper previously wrote the knowledge
% matrix two ways -- M(W,f)(x) in the abstract/intro/setup and M(theta,x) in the
% math core -- which was tolerable while the network function was f(theta,.),
% because the theta matched. With the network function now Psi(W,f), the (W,f)
% form is canonical and the theta-indexed form reads as a leftover. These macros
% carry (W,f) into the theorem statements and proofs.
%   \KMwf        M(W,f)          \KMwfx        M(W,f)(x)
%   \KMarg{W'}{x}  M(W',f)(x)    -- for a second/transformed weight collection
%   \Jwf         J(W,f)          \cwf          c(W,f)
\newcommand{\KMwf}{\M(W\!,f)}
\newcommand{\KMwfx}{\M(W\!,f)(x)}
\newcommand{\KMarg}[2]{\M(#1\!,f)(#2)}
\newcommand{\Jwf}{J(W\!,f)}
\newcommand{\cwf}{c(W\!,f)}
). Edit document-wide formatting here.
% ============================================================================

% --- Packages ---
\usepackage[utf8]{inputenc}
\usepackage[T1]{fontenc}
\usepackage{lmodern}
\usepackage[margin=1in]{geometry}
\usepackage{amsmath,amssymb,amsthm}
\usepackage{booktabs}
\usepackage{graphicx}
\usepackage{enumitem}
\usepackage{xcolor}
\usepackage[colorlinks=true,linkcolor=blue!70!black,citecolor=green!50!black,urlcolor=blue!70!black]{hyperref}
\usepackage[round]{natbib}
\usepackage{tikz}
\usetikzlibrary{arrows.meta,decorations.pathreplacing,calc,positioning}
\usepackage[most]{tcolorbox}

% --- Theorem environments ---
\newtheorem{theorem}{Theorem}
\newtheorem{proposition}[theorem]{Proposition}
\newtheorem{definition}[theorem]{Definition}
\newtheorem{remark}[theorem]{Remark}
\newtheorem{lemma}[theorem]{Lemma}
\newtheorem{corollary}[theorem]{Corollary}
\newtheorem{conjecture}[theorem]{Conjecture}
\numberwithin{theorem}{section}

% --- Notation ---
\newcommand{\R}{\mathbb{R}}
\newcommand{\M}{\mathrm{M}}
\newcommand{\Mfx}{M(W{,}f)(x)}
\newcommand{\Mfxp}{M(W{,}f)(x')}
\newcommand{\Weff}{W_{\text{eff}}}
\newcommand{\beff}{b_{\text{eff}}}
\newcommand{\KM}{\textsc{KM}}

% Placeholder macro for empirical numbers wired in by wire_paper_results.py.
% Renders the argument verbatim in the PDF until substituted.
\newcommand{\TODO}[1]{\textcolor{red}{\textbf{[TODO:}\ #1\textbf{]}}}

% ===== REVIEW-ONLY WARNING BOXES (remove every \WARNING{...} before submission) =====
% Red-framed [WARNING] box with a detailed explanation, visible in the compiled PDF.
\newcommand{\WARNING}[1]{%
  \par\smallskip\noindent
  \fcolorbox{red}{red!6}{%
    \begin{minipage}{0.93\linewidth}%
      \small\textcolor{red}{\textbf{[WARNING --- review note; remove before submission]}}\par\smallskip
      #1%
    \end{minipage}}\par\smallskip}

% ===== INLINE CLARIFICATIONS (added 2026-09-07) =====
% Blue inline text answering reader questions: definitions of notation and
% terminology given at the point of first use, rather than in a glossary.
%   \clarify{...}     inline blue prose (a sentence or a short paragraph)
%   \clarifyterm{t}{d} bolds the term being defined, then the blue definition
% Both are plain \textcolor wrappers -- they add no vertical space and can be
% used inside theorem statements, captions, and tables.
\definecolor{clarifyblue}{RGB}{0,70,190}
\newcommand{\clarify}[1]{\textcolor{clarifyblue}{#1}}
\newcommand{\clarifyterm}[2]{\textcolor{clarifyblue}{\emph{#1}: #2}}
% Display-style clarification: an indented blue block for multi-sentence notes.
\newenvironment{clarifybox}
  {\par\smallskip\begingroup\color{clarifyblue}\small\leftskip=1.2em\rightskip=1.2em\noindent\ignorespaces}
  {\par\endgroup\smallskip}

% ===== NETWORK FUNCTION: Psi(W,f) (notation fix, 2026-09-07) =====
% The paper previously used f for BOTH the activation function (as in the
% knowledge matrix M(W,f), where f IS the activation) and the network function
% (as in f(theta,x)). That collision is resolved by writing the network function
% in the quiver-representation notation of Armenta & Jodoin: Psi(W,f), the
% function realised by weight collection W and activation f. From here on
%   f      = the activation function
%   Psi(W,f) = the network function it realises
\newcommand{\Net}{\Psi(W\!,f)}          % the network function itself
\newcommand{\Netx}{\Psi(W\!,f)(x)}      % evaluated at x
\newcommand{\Nety}{\Psi(W\!,f)(y)}
\newcommand{\Netxp}{\Psi(W\!,f)(x')}
\newcommand{\Netarg}[1]{\Psi(W\!,f)(#1)}

% ===== KNOWLEDGE MATRIX indexed by (W,f), not by theta (2026-09-07) =====
% Companion to the \Net macros above. The paper previously wrote the knowledge
% matrix two ways -- M(W,f)(x) in the abstract/intro/setup and M(theta,x) in the
% math core -- which was tolerable while the network function was f(theta,.),
% because the theta matched. With the network function now Psi(W,f), the (W,f)
% form is canonical and the theta-indexed form reads as a leftover. These macros
% carry (W,f) into the theorem statements and proofs.
%   \KMwf        M(W,f)          \KMwfx        M(W,f)(x)
%   \KMarg{W'}{x}  M(W',f)(x)    -- for a second/transformed weight collection
%   \Jwf         J(W,f)          \cwf          c(W,f)
\newcommand{\KMwf}{\M(W\!,f)}
\newcommand{\KMwfx}{\M(W\!,f)(x)}
\newcommand{\KMarg}[2]{\M(#1\!,f)(#2)}
\newcommand{\Jwf}{J(W\!,f)}
\newcommand{\cwf}{c(W\!,f)}
). Edit document-wide formatting here.
% ============================================================================

% --- Packages ---
\usepackage[utf8]{inputenc}
\usepackage[T1]{fontenc}
\usepackage{lmodern}
\usepackage[margin=1in]{geometry}
\usepackage{amsmath,amssymb,amsthm}
\usepackage{booktabs}
\usepackage{graphicx}
\usepackage{enumitem}
\usepackage{xcolor}
\usepackage[colorlinks=true,linkcolor=blue!70!black,citecolor=green!50!black,urlcolor=blue!70!black]{hyperref}
\usepackage[round]{natbib}
\usepackage{tikz}
\usetikzlibrary{arrows.meta,decorations.pathreplacing,calc,positioning}
\usepackage[most]{tcolorbox}

% --- Theorem environments ---
\newtheorem{theorem}{Theorem}
\newtheorem{proposition}[theorem]{Proposition}
\newtheorem{definition}[theorem]{Definition}
\newtheorem{remark}[theorem]{Remark}
\newtheorem{lemma}[theorem]{Lemma}
\newtheorem{corollary}[theorem]{Corollary}
\newtheorem{conjecture}[theorem]{Conjecture}
\numberwithin{theorem}{section}

% --- Notation ---
\newcommand{\R}{\mathbb{R}}
\newcommand{\M}{\mathrm{M}}
\newcommand{\Mfx}{M(W{,}f)(x)}
\newcommand{\Mfxp}{M(W{,}f)(x')}
\newcommand{\Weff}{W_{\text{eff}}}
\newcommand{\beff}{b_{\text{eff}}}
\newcommand{\KM}{\textsc{KM}}

% Placeholder macro for empirical numbers wired in by wire_paper_results.py.
% Renders the argument verbatim in the PDF until substituted.
\newcommand{\TODO}[1]{\textcolor{red}{\textbf{[TODO:}\ #1\textbf{]}}}

% ===== REVIEW-ONLY WARNING BOXES (remove every \WARNING{...} before submission) =====
% Red-framed [WARNING] box with a detailed explanation, visible in the compiled PDF.
\newcommand{\WARNING}[1]{%
  \par\smallskip\noindent
  \fcolorbox{red}{red!6}{%
    \begin{minipage}{0.93\linewidth}%
      \small\textcolor{red}{\textbf{[WARNING --- review note; remove before submission]}}\par\smallskip
      #1%
    \end{minipage}}\par\smallskip}

% ===== INLINE CLARIFICATIONS (added 2026-09-07) =====
% Blue inline text answering reader questions: definitions of notation and
% terminology given at the point of first use, rather than in a glossary.
%   \clarify{...}     inline blue prose (a sentence or a short paragraph)
%   \clarifyterm{t}{d} bolds the term being defined, then the blue definition
% Both are plain \textcolor wrappers -- they add no vertical space and can be
% used inside theorem statements, captions, and tables.
\definecolor{clarifyblue}{RGB}{0,70,190}
\newcommand{\clarify}[1]{\textcolor{clarifyblue}{#1}}
\newcommand{\clarifyterm}[2]{\textcolor{clarifyblue}{\emph{#1}: #2}}
% Display-style clarification: an indented blue block for multi-sentence notes.
\newenvironment{clarifybox}
  {\par\smallskip\begingroup\color{clarifyblue}\small\leftskip=1.2em\rightskip=1.2em\noindent\ignorespaces}
  {\par\endgroup\smallskip}

% ===== NETWORK FUNCTION: Psi(W,f) (notation fix, 2026-09-07) =====
% The paper previously used f for BOTH the activation function (as in the
% knowledge matrix M(W,f), where f IS the activation) and the network function
% (as in f(theta,x)). That collision is resolved by writing the network function
% in the quiver-representation notation of Armenta & Jodoin: Psi(W,f), the
% function realised by weight collection W and activation f. From here on
%   f      = the activation function
%   Psi(W,f) = the network function it realises
\newcommand{\Net}{\Psi(W\!,f)}          % the network function itself
\newcommand{\Netx}{\Psi(W\!,f)(x)}      % evaluated at x
\newcommand{\Nety}{\Psi(W\!,f)(y)}
\newcommand{\Netxp}{\Psi(W\!,f)(x')}
\newcommand{\Netarg}[1]{\Psi(W\!,f)(#1)}

% ===== KNOWLEDGE MATRIX indexed by (W,f), not by theta (2026-09-07) =====
% Companion to the \Net macros above. The paper previously wrote the knowledge
% matrix two ways -- M(W,f)(x) in the abstract/intro/setup and M(theta,x) in the
% math core -- which was tolerable while the network function was f(theta,.),
% because the theta matched. With the network function now Psi(W,f), the (W,f)
% form is canonical and the theta-indexed form reads as a leftover. These macros
% carry (W,f) into the theorem statements and proofs.
%   \KMwf        M(W,f)          \KMwfx        M(W,f)(x)
%   \KMarg{W'}{x}  M(W',f)(x)    -- for a second/transformed weight collection
%   \Jwf         J(W,f)          \cwf          c(W,f)
\newcommand{\KMwf}{\M(W\!,f)}
\newcommand{\KMwfx}{\M(W\!,f)(x)}
\newcommand{\KMarg}[2]{\M(#1\!,f)(#2)}
\newcommand{\Jwf}{J(W\!,f)}
\newcommand{\cwf}{c(W\!,f)}
). Edit document-wide formatting here.
% ============================================================================

% --- Packages ---
\usepackage[utf8]{inputenc}
\usepackage[T1]{fontenc}
\usepackage{lmodern}
\usepackage[margin=1in]{geometry}
\usepackage{amsmath,amssymb,amsthm}
\usepackage{booktabs}
\usepackage{graphicx}
\usepackage{enumitem}
\usepackage{xcolor}
\usepackage[colorlinks=true,linkcolor=blue!70!black,citecolor=green!50!black,urlcolor=blue!70!black]{hyperref}
\usepackage[round]{natbib}
\usepackage{tikz}
\usetikzlibrary{arrows.meta,decorations.pathreplacing,calc,positioning}
\usepackage[most]{tcolorbox}

% --- Theorem environments ---
\newtheorem{theorem}{Theorem}
\newtheorem{proposition}[theorem]{Proposition}
\newtheorem{definition}[theorem]{Definition}
\newtheorem{remark}[theorem]{Remark}
\newtheorem{lemma}[theorem]{Lemma}
\newtheorem{corollary}[theorem]{Corollary}
\newtheorem{conjecture}[theorem]{Conjecture}
\numberwithin{theorem}{section}

% --- Notation ---
\newcommand{\R}{\mathbb{R}}
\newcommand{\M}{\mathrm{M}}
\newcommand{\Mfx}{M(W{,}f)(x)}
\newcommand{\Mfxp}{M(W{,}f)(x')}
\newcommand{\Weff}{W_{\text{eff}}}
\newcommand{\beff}{b_{\text{eff}}}
\newcommand{\KM}{\textsc{KM}}

% Placeholder macro for empirical numbers wired in by wire_paper_results.py.
% Renders the argument verbatim in the PDF until substituted.
\newcommand{\TODO}[1]{\textcolor{red}{\textbf{[TODO:}\ #1\textbf{]}}}

% ===== REVIEW-ONLY WARNING BOXES (remove every \WARNING{...} before submission) =====
% Red-framed [WARNING] box with a detailed explanation, visible in the compiled PDF.
\newcommand{\WARNING}[1]{%
  \par\smallskip\noindent
  \fcolorbox{red}{red!6}{%
    \begin{minipage}{0.93\linewidth}%
      \small\textcolor{red}{\textbf{[WARNING --- review note; remove before submission]}}\par\smallskip
      #1%
    \end{minipage}}\par\smallskip}

% ===== INLINE CLARIFICATIONS (added 2026-09-07) =====
% Blue inline text answering reader questions: definitions of notation and
% terminology given at the point of first use, rather than in a glossary.
%   \clarify{...}     inline blue prose (a sentence or a short paragraph)
%   \clarifyterm{t}{d} bolds the term being defined, then the blue definition
% Both are plain \textcolor wrappers -- they add no vertical space and can be
% used inside theorem statements, captions, and tables.
\definecolor{clarifyblue}{RGB}{0,70,190}
\newcommand{\clarify}[1]{\textcolor{clarifyblue}{#1}}
\newcommand{\clarifyterm}[2]{\textcolor{clarifyblue}{\emph{#1}: #2}}
% Display-style clarification: an indented blue block for multi-sentence notes.
\newenvironment{clarifybox}
  {\par\smallskip\begingroup\color{clarifyblue}\small\leftskip=1.2em\rightskip=1.2em\noindent\ignorespaces}
  {\par\endgroup\smallskip}

% ===== NETWORK FUNCTION: Psi(W,f) (notation fix, 2026-09-07) =====
% The paper previously used f for BOTH the activation function (as in the
% knowledge matrix M(W,f), where f IS the activation) and the network function
% (as in f(theta,x)). That collision is resolved by writing the network function
% in the quiver-representation notation of Armenta & Jodoin: Psi(W,f), the
% function realised by weight collection W and activation f. From here on
%   f      = the activation function
%   Psi(W,f) = the network function it realises
\newcommand{\Net}{\Psi(W\!,f)}          % the network function itself
\newcommand{\Netx}{\Psi(W\!,f)(x)}      % evaluated at x
\newcommand{\Nety}{\Psi(W\!,f)(y)}
\newcommand{\Netxp}{\Psi(W\!,f)(x')}
\newcommand{\Netarg}[1]{\Psi(W\!,f)(#1)}

% ===== KNOWLEDGE MATRIX indexed by (W,f), not by theta (2026-09-07) =====
% Companion to the \Net macros above. The paper previously wrote the knowledge
% matrix two ways -- M(W,f)(x) in the abstract/intro/setup and M(theta,x) in the
% math core -- which was tolerable while the network function was f(theta,.),
% because the theta matched. With the network function now Psi(W,f), the (W,f)
% form is canonical and the theta-indexed form reads as a leftover. These macros
% carry (W,f) into the theorem statements and proofs.
%   \KMwf        M(W,f)          \KMwfx        M(W,f)(x)
%   \KMarg{W'}{x}  M(W',f)(x)    -- for a second/transformed weight collection
%   \Jwf         J(W,f)          \cwf          c(W,f)
\newcommand{\KMwf}{\M(W\!,f)}
\newcommand{\KMwfx}{\M(W\!,f)(x)}
\newcommand{\KMarg}[2]{\M(#1\!,f)(#2)}
\newcommand{\Jwf}{J(W\!,f)}
\newcommand{\cwf}{c(W\!,f)}
